\documentclass[11pt,a4paper]{article}

% --- Packages ---
\usepackage[utf8]{inputenc}
\usepackage[T1]{fontenc}
\usepackage{lmodern}
\usepackage[margin=2.5cm,headheight=14pt]{geometry}
\usepackage{amsmath,amssymb,amsthm}
\usepackage{mathtools}
\usepackage{enumitem}
\usepackage{booktabs}
\usepackage{graphicx}
\usepackage{hyperref}
\usepackage{cleveref}
\usepackage{xcolor}
\usepackage{tcolorbox}
\usepackage{fancyhdr}
\usepackage{titlesec}
\usepackage{caption}

% --- Colors ---
\definecolor{olsblue}{HTML}{2E5EA8}
\definecolor{truthgold}{HTML}{D4A017}
\definecolor{biasred}{HTML}{C0392B}
\definecolor{repairgreen}{HTML}{27AE60}

% --- Theorem environments ---
\theoremstyle{definition}
\newtheorem{definition}{Definition}[section]
\newtheorem{example}{Example}[section]
\theoremstyle{plain}
\newtheorem{theorem}{Theorem}[section]
\newtheorem{proposition}{Proposition}[section]
\newtheorem{lemma}{Lemma}[section]
\newtheorem{corollary}{Corollary}[section]
\theoremstyle{remark}
\newtheorem{remark}{Remark}[section]

% --- Coloured boxes ---
\tcbuselibrary{breakable}
\newtcolorbox{keyinsight}[1][]{colback=olsblue!5, colframe=olsblue!80!black,
  fonttitle=\bfseries, title={Key Insight}, breakable, #1}
\newtcolorbox{warningbox}[1][]{colback=biasred!5, colframe=biasred!80!black,
  fonttitle=\bfseries, title={Warning}, breakable, #1}
\newtcolorbox{repairbox}[1][]{colback=repairgreen!5, colframe=repairgreen!80!black,
  fonttitle=\bfseries, title={The Repair}, breakable, #1}

% --- New pedagogical boxes ---
\newtcolorbox{toystory}[1][]{colback=truthgold!5, colframe=truthgold!80!black,
  fonttitle=\bfseries, title={Toy Story}, breakable, #1}
\newtcolorbox{numermark}[1][]{colback=olsblue!3, colframe=olsblue!50!black,
  fonttitle=\bfseries, title={Numerical Example}, breakable, #1}
\newtcolorbox{vizspec}[1][]{colback=repairgreen!3, colframe=repairgreen!50!black,
  fonttitle=\bfseries, title={Visualization}, breakable, #1}

% --- Header / Footer ---
\pagestyle{fancy}
\fancyhf{}
\lhead{\textsc{Regression Autopsy}}
\rhead{\textsc{Lecture Notes}}
\cfoot{\thepage}
\renewcommand{\headrulewidth}{0.4pt}

% --- Section formatting ---
\titleformat{\section}{\Large\bfseries\color{olsblue!80!black}}{\thesection}{1em}{}
\titleformat{\subsection}{\large\bfseries\color{olsblue!60!black}}{\thesubsection}{1em}{}

% --- Hyperref setup ---
\hypersetup{
  colorlinks=true,
  linkcolor=olsblue,
  citecolor=olsblue,
  urlcolor=olsblue!70!black,
}

% --- Custom commands ---
\newcommand{\E}{\mathbb{E}}
\newcommand{\Var}{\operatorname{Var}}
\newcommand{\Cov}{\operatorname{Cov}}
\newcommand{\Bias}{\operatorname{Bias}}
\newcommand{\MSE}{\operatorname{MSE}}
\newcommand{\se}{\operatorname{SE}}
\newcommand{\plim}{\operatorname{plim}}
\newcommand{\bhat}{\hat{\beta}}
\newcommand{\bols}{\hat{\beta}_{\text{OLS}}}
\newcommand{\biv}{\hat{\beta}_{\text{IV}}}

\begin{document}

% ============================================================
% TITLE PAGE
% ============================================================
\begin{titlepage}
\centering
\vspace*{3cm}
{\Huge\bfseries Regression Autopsy\par}
\vspace{0.5cm}
{\LARGE\color{olsblue!80!black}\itshape Eight Ways Your Model Is Lying to You\par}
\vspace{2cm}
{\Large Lecture Notes\par}
\vspace{1cm}
{\large Statistical and Machine Learning Lab\par}
\vspace{0.5cm}
{\large Spring 2026\par}
\vfill
\rule{\textwidth}{0.4pt}\\[0.3cm]
{\small These notes accompany the interactive Jupyter notebook series.
All theorems stated here are validated by Monte Carlo simulation in the corresponding notebooks.}
\end{titlepage}

\tableofcontents
\newpage

% ============================================================
% 1. INTRODUCTION
% ============================================================
\section{Introduction}\label{sec:intro}

Regression is the most widely used tool in quantitative data analysis.
It produces point estimates, standard errors, confidence intervals, $p$-values, and $R^2$ statistics---all of which look authoritative.
In most observational analyses, at least one of these numbers is wrong.

These lecture notes formalise the mathematical foundations behind the eight core failure modes of ordinary least squares (OLS) regression.
Each section corresponds to one interactive notebook in the accompanying lab series.
The notes are designed to be self-contained: every claim is stated as a formal result and, in the lab, confirmed by Monte Carlo simulation.

Each section follows a consistent pattern: we begin with an everyday story that makes the problem visceral, then state the rigorous mathematics, then work through a tiny numerical example you can verify on paper, and finally describe an interactive visualization that lets you play with the result.

\subsection{The NovaMart Scenario}

Throughout the course, we work with a single fictional company---\textbf{NovaMart}, an online retailer.
The variables of interest are:
\begin{itemize}[nosep]
  \item $Y$: monthly revenue (thousands of dollars),
  \item $X_1$: advertising spend (thousands of dollars),
  \item $X_2$: number of sales staff,
  \item $X_3$: customer satisfaction score,
  \item $U$: unobserved market demand (hidden from the analyst).
\end{itemize}

Using a single recurring dataset creates continuity: when a coefficient changes from one notebook to the next, the student remembers the old value.

% ============================================================
% 2. THE DATA GENERATING PROCESS
% ============================================================
\section{The Data Generating Process}\label{sec:dgp}

\begin{definition}[NovaMart DGP]\label{def:dgp}
The data are generated according to the following structural equations:
\begin{align}
  U &\sim \mathcal{N}(0,1), \label{eq:U}\\
  X_1 &= \gamma_1 U + \eta_1, \quad \eta_1 \sim \mathcal{N}(0,\sigma_\eta^2), \label{eq:X1}\\
  X_2 &= \gamma_2 U + \eta_2, \quad \eta_2 \sim \mathcal{N}(0,\sigma_\eta^2), \label{eq:X2}\\
  \varepsilon &\sim \mathcal{N}\!\left(0,\; h \cdot |X_1|\right), \label{eq:eps}\\
  Y &= \beta_0 + \beta_1 \log(X_1) + \beta_2 X_2 + \beta_U U + \varepsilon, \label{eq:Y}\\
  X_3 &= \alpha \, Y + \eta_3, \quad \eta_3 \sim \mathcal{N}(0,\sigma_\eta^2). \label{eq:X3}
\end{align}
\end{definition}

If these structural equations look intimidating, don't worry. Each section ahead will isolate one piece and make it tangible through stories and examples before we prove anything.

\begin{remark}[Probability limit]
Throughout these notes, $\plim$ denotes the \emph{probability limit}---the value an estimator converges to as the sample size grows to infinity. Think of it as ``what the estimator settles down to with unlimited data.''
\end{remark}

The default parameter values are:
\begin{center}
\begin{tabular}{@{}lll@{}}
\toprule
\textbf{Parameter} & \textbf{Symbol} & \textbf{Default} \\
\midrule
Intercept & $\beta_0$ & 50 \\
Ad spend effect & $\beta_1$ & 8 \\
Staff effect & $\beta_2$ & 3 \\
Demand effect & $\beta_U$ & 2 \\
Endogeneity strength & $\gamma_1$ & 20 \\
Staff--demand loading & $\gamma_2$ & 5 \\
Heteroscedasticity & $h$ & 0.5 \\
Bad-control strength & $\alpha$ & 0.1 \\
Noise scale & $\sigma_\eta$ & 1 \\
\bottomrule
\end{tabular}
\end{center}

\begin{remark}[Causal ordering]
The generation order is critical: first $U$, then $(X_1, X_2)$, then $Y$, then $X_3$.
The variable $X_3$ is \emph{post-treatment}: it is caused by $Y$, not a cause of $Y$.
It never appears in the true DGP for~$Y$.
Including $X_3$ as a control variable introduces collider bias (see \Cref{sec:causation}).
\end{remark}

\subsection{Pathologies Embedded in the DGP}

The single DGP simultaneously contains all of the following pathologies:
\begin{enumerate}[nosep]
  \item \textbf{Omitted variable bias:} $U$ drives both $X_1$ and $Y$.
        Omitting $U$ biases $\bhat_1$ upward.
  \item \textbf{Heteroscedasticity:} $\Var(\varepsilon \mid X_1) = h \cdot |X_1|$;
        larger campaigns have more variable returns.
  \item \textbf{Inflated significance:} with large $n$, trivially small effects become ``significant.''
  \item \textbf{Misspecification:} the true relationship is $\log(X_1)$, not $X_1$.
  \item \textbf{Overfitting:} adding polynomial terms inflates training $R^2$ but destroys test $R^2$.
  \item \textbf{Collider bias / bad control:} $X_3$ is caused by $Y$;
        conditioning on it opens a spurious path.
\end{enumerate}

Each notebook isolates and examines one pathology while the others remain present in the background.

% ============================================================
% 3. OMITTED VARIABLE BIAS (NB 1)
% ============================================================
\section{Why Your Coefficient Is Wrong: Omitted Variable Bias}\label{sec:ovb}

\subsection{Setup}

\begin{toystory}[title={The Sunscreen Murder Mystery}]
Every summer, ice cream sales go up and so do drowning deaths. If you just look at the data, it seems like ice cream \emph{causes} drownings. But you forgot about the hot weather---when it's hot, people buy more ice cream AND swim more. The heat is doing the real work behind the scenes, but because you ignored it, its effect got smeared onto ice cream, making ice cream look dangerous.

This is exactly what the OVB formula captures: the bias on ice cream's coefficient equals the true effect of temperature on drownings times the correlation between temperature and ice cream sales. Leave out a variable that matters and is correlated with what you included, and its ghost haunts your estimate.
\end{toystory}

Suppose the true model is
\begin{equation}\label{eq:long}
  Y = \beta_0 + \beta_1 X_1 + \beta_2 X_2 + \varepsilon, \qquad \E[\varepsilon \mid X_1, X_2] = 0,
\end{equation}
but the analyst estimates the \emph{short} regression
\begin{equation}\label{eq:short}
  Y = \tilde\beta_0 + \tilde\beta_1 X_1 + \tilde\varepsilon.
\end{equation}

\begin{theorem}[Omitted Variable Bias Formula]\label{thm:ovb}
Let $\bhat_1^{\text{short}}$ denote the OLS estimator of $\tilde\beta_1$ from the short regression~\eqref{eq:short}, and let $\beta_1$ be the true coefficient in~\eqref{eq:long}.  Define the \emph{auxiliary slope}
\[
  \delta_1 \;=\; \frac{\Cov(X_1, X_2)}{\Var(X_1)}.
\]
Then
\begin{equation}\label{eq:ovb}
  \E[\bhat_1^{\emph{short}}] \;=\; \beta_1 + \beta_2 \, \delta_1.
\end{equation}
The term $\beta_2 \delta_1$ is the \textbf{omitted variable bias}.
\end{theorem}

\begin{proof}
Write the OLS estimator in population form.  From the short regression,
\[
  \plim\, \bhat_1^{\text{short}}
  = \frac{\Cov(X_1, Y)}{\Var(X_1)}
  = \frac{\Cov(X_1,\, \beta_0 + \beta_1 X_1 + \beta_2 X_2 + \varepsilon)}{\Var(X_1)}.
\]
Since $\E[\varepsilon \mid X_1, X_2] = 0$ implies $\Cov(X_1, \varepsilon) = 0$, we obtain
\[
  \plim\, \bhat_1^{\text{short}}
  = \beta_1 + \beta_2 \cdot \frac{\Cov(X_1, X_2)}{\Var(X_1)}
  = \beta_1 + \beta_2 \, \delta_1. \qedhere
\]
\end{proof}

\begin{numermark}[title={Studying, Sleeping, and Exam Scores}]
Five students have data on hours studied ($X_1$), hours slept ($X_2$), and exam scores ($Y$). The true model is $Y = 60 + 5X_1 + 3X_2$ (no noise).

\begin{center}
\begin{tabular}{@{}lccc@{}}
\toprule
Student & $X_1$ (studied) & $X_2$ (slept) & $Y$ (score) \\
\midrule
A & 1 & 8 & 89 \\
B & 2 & 6 & 88 \\
C & 3 & 7 & 96 \\
D & 4 & 4 & 92 \\
E & 5 & 5 & 100 \\
\bottomrule
\end{tabular}
\end{center}

\textbf{Full regression} (include both $X_1$ and $X_2$): recovers $\bhat_1 = 5$, $\bhat_2 = 3$.

\textbf{Short regression} (omit sleep): gives $\bhat_1^{\text{short}} = 2.6$.

\textbf{Auxiliary regression} ($X_2$ on $X_1$): slope $\delta_1 = -0.8$ (each extra hour of study costs 0.8 hours of sleep).

\textbf{Verify the OVB formula:}
\[
  \bhat_1^{\text{short}} = \beta_1 + \beta_2 \cdot \delta_1 = 5 + 3 \times (-0.8) = 5 - 2.4 = 2.6. \quad \checkmark
\]

Omitting sleep shrank the study coefficient from 5 to 2.6---each hour of study costs sleep, and sleep helps scores, so some of study's real effect gets ``charged'' to the missing sleep variable.
\end{numermark}

\begin{vizspec}[title={``The Ghost in the Regression'' --- Interactive OVB Explorer}]
\textbf{Try it:} A scatterplot of $X_1$ vs.\ $Y$ shows two regression lines---the full model (gold) and the short model (red). Toggle a switch to include/exclude $X_2$ and watch the slope shift in real time.

\textbf{Key interactions:}
\begin{itemize}[nosep]
  \item \emph{Correlation slider} ($-0.95$ to $+0.95$): controls how strongly $X_1$ and $X_2$ are related. At zero correlation, both lines coincide---omitting an uncorrelated variable causes no bias.
  \item \emph{Effect-size slider} ($\beta_2$ from $-5$ to $+10$): when the omitted variable doesn't matter ($\beta_2 = 0$), there's no bias regardless of correlation.
\end{itemize}
The formula panel updates live, showing how $\beta_2 \cdot \delta_1$ produces the exact gap you see.
\end{vizspec}

\begin{keyinsight}
The bias has a \emph{sign} and a \emph{direction}.
In the NovaMart DGP with default parameters, $\beta_2 = 3$ and $\delta_1 > 0$ (because both $X_1$ and $X_2$ are driven by $U$), so the bias is \emph{positive}: the short regression overestimates the effect of ad spend.
The OVB formula predicts $\E[\bhat_1^{\text{short}}] \approx 2.3 + 2.0 \times 1.2 = 4.7$, which is confirmed by Monte Carlo simulation (mean $\approx 4.698$ over 5{,}000 replications).
\end{keyinsight}

\subsection{When OVB Cannot Be Signed}

The formula~\eqref{eq:ovb} generalises to multiple omitted variables.
If $k$ variables are omitted, the bias on $\bhat_1$ is a weighted sum of terms $\beta_j \delta_{1j}$, where each $\delta_{1j}$ is the partial regression coefficient of $X_j$ on $X_1$ after partialling out included regressors.
In general, the sign of the total bias is indeterminate without knowledge of all omitted variables.

\subsection{The Proxy Variable Approach}

\begin{repairbox}
If the true confounder $U$ is unobservable, one can include a \emph{proxy} $\tilde{U}$ that is correlated with $U$.
If $\tilde{U}$ explains a fraction $q \in [0,1]$ of the variation in $U$, the residual bias is approximately $(1-q)$ times the original bias.
At $q = 1$ (perfect proxy), the bias vanishes; at $q = 0$, the full bias remains.
\end{repairbox}

\begin{warningbox}
Including $X_3$ (customer satisfaction) as a control makes things \emph{worse}, not better.
Because $X_3$ is caused by $Y$, conditioning on $X_3$ opens a collider path and introduces additional bias.
\end{warningbox}

% ============================================================
% 4. HETEROSCEDASTICITY AND STANDARD ERRORS (NB 2)
% ============================================================
\section{Why Your Uncertainty Is Wrong: Heteroscedasticity}\label{sec:hetero}

\subsection{The Classical Variance Formula}

\begin{toystory}[title={The Bathroom Scale That Lies About Its Own Reliability}]
You have a bathroom scale that gives the right weight \emph{on average}---step on it 100 times and the average is your true weight. But when you weigh a kitten, the readings barely wobble. When you weigh a suitcase full of books, the readings scatter all over the place. If you pretend the wobble is the same for everything, you'll think you know the suitcase's weight much more precisely than you actually do. The scale isn't biased---it's centered on the truth---but your confidence about ``how close am I?'' is completely wrong.

This is what the sandwich variance formula fixes: when the noise differs across observations, the simple formula for uncertainty is wrong, and you need the corrected version to get honest error bars.
\end{toystory}

Under the standard linear model $Y = X\beta + \varepsilon$ with $\E[\varepsilon\varepsilon^\top \mid X] = \sigma^2 I_n$, the variance of the OLS estimator is
\begin{equation}\label{eq:var-homo}
  \Var(\bols \mid X) = \sigma^2 (X^\top X)^{-1}.
\end{equation}

\begin{definition}[Heteroscedasticity]
The error term is \emph{heteroscedastic} if $\Var(\varepsilon_i \mid X_i) = \sigma_i^2$ with $\sigma_i^2$ not constant across $i$.
In the NovaMart DGP, $\sigma_i^2 = h \cdot |X_{1,i}|$: larger ad campaigns have more unpredictable returns.
\end{definition}

\begin{theorem}[Variance Under Heteroscedasticity]\label{thm:sandwich}
When $\E[\varepsilon\varepsilon^\top \mid X] = \Omega \neq \sigma^2 I_n$, the OLS estimator $\bols$ remains \emph{unbiased} but its variance becomes
\begin{equation}\label{eq:var-hetero}
  \Var(\bols \mid X) = (X^\top X)^{-1} X^\top \Omega\, X\, (X^\top X)^{-1}.
\end{equation}
This is the \textbf{sandwich formula}.
The classical formula~\eqref{eq:var-homo} underestimates or overestimates the true variance, depending on the correlation structure between $X$ and the error variances.
\end{theorem}

\begin{numermark}[title={Five Companies, One Funnel}]
We predict revenue ($Y$, in \$K) from ad spend ($X$, in \$K) for five companies:

\begin{center}
\begin{tabular}{@{}lccccc@{}}
\toprule
Company & $X$ & $Y$ & $\hat{Y} = 2.2X$ & Residual $\hat\varepsilon$ & $\hat\varepsilon^2$ \\
\midrule
A & 1 & 3 & 2.2 & $+0.8$ & 0.64 \\
B & 2 & 5 & 4.4 & $+0.6$ & 0.36 \\
C & 3 & 7 & 6.6 & $+0.4$ & 0.16 \\
D & 4 & 3 & 8.8 & $-5.8$ & 33.64 \\
E & 5 & 15 & 11.0 & $+4.0$ & 16.00 \\
\bottomrule
\end{tabular}
\end{center}

Low-$X$ companies (A, B) have average $\hat\varepsilon^2 = 0.50$. High-$X$ companies (D, E) have average $\hat\varepsilon^2 = 24.82$---\textbf{fifty times larger}. The classical SE of $1.30$ treats all five companies as equally informative. But the big spenders are wildly noisier, so the classical confidence interval is a lie.
\end{numermark}

\begin{vizspec}[title={``The Funnel of False Confidence'' --- Coverage Simulator}]
\textbf{Try it:} A scatterplot shows data fanning out (the funnel). Two confidence bands overlay the regression line: classical (constant width, blue) and robust (varying width, green).

\textbf{Key interaction:} Click ``Draw 100 new samples.'' Watch 100 confidence intervals cascade in, coloring green if they contain the truth and red if they miss. With heteroscedasticity, classical CIs achieve only $\approx$82\% coverage while robust CIs hit the promised 95\%.

\textbf{Break it:} Set sample size to $n = 20$ and watch the robust SEs \emph{perform worse} than classical---the repair has a minimum sample-size cost.
\end{vizspec}

\subsection{Consequences for Inference}

If the classical standard errors are used under heteroscedasticity:
\begin{itemize}[nosep]
  \item Confidence intervals have incorrect coverage.
        In the NovaMart simulation, the nominal 95\% CI achieves only $\approx$82\% coverage.
  \item $t$-statistics and $p$-values are unreliable.
  \item The coefficient estimates themselves are \emph{not} biased---only the uncertainty is wrong.
\end{itemize}

\subsection{The Breusch--Pagan Test}

\begin{definition}[Breusch--Pagan Test]
Regress the squared OLS residuals $\hat\varepsilon_i^2$ on the regressors $X$.
Under $H_0$: homoscedasticity, the test statistic $nR^2$ from this auxiliary regression follows $\chi^2_p$ asymptotically, where $p$ is the number of regressors (excluding the constant).
\end{definition}

\subsection{Robust Standard Errors}

\begin{repairbox}
The \textbf{White--Huber--Eicker} (``sandwich'') standard errors are computed as
\[
  \widehat{\Var}(\bols) = (X^\top X)^{-1} \left(\sum_{i=1}^n \hat\varepsilon_i^2\, x_i x_i^\top\right) (X^\top X)^{-1},
\]
where $\hat\varepsilon_i$ are the OLS residuals.  These standard errors are consistent under heteroscedasticity of unknown form.
With robust SEs, the 95\% CI recovers $\approx$95\% coverage.
\end{repairbox}

\begin{warningbox}
Robust standard errors perform poorly in small samples.
When $n = 20$, the sandwich estimator is more biased than the classical estimator.
The crossover point depends on the severity of heteroscedasticity, but as a rough guide, $n \geq 50$ is typically needed for reliable sandwich SEs.
\end{warningbox}

% ============================================================
% 5. SIGNIFICANCE VS EFFECT SIZE (NB 3)
% ============================================================
\section{Why Your Significance Is Wrong: Effect Size and Power}\label{sec:significance}

\subsection{The Anatomy of a $t$-Statistic}

\begin{toystory}[title={The Coin That's Almost Fair}]
Imagine a coin that lands heads 50.1\% of the time instead of exactly 50\%. Flip it 20 times and you'd never notice---it looks perfectly fair. Flip it 10 million times and you can prove beyond any doubt that it's rigged. But so what? It's still basically fair---no casino would care, no bet would change. You ``detected'' a difference so tiny it means nothing in real life, just because you had an ocean of data. Statistical significance told you ``it's not exactly 50\%,'' but it never told you ``and you should care.''

This is exactly what the $t$-statistic scaling with $\sqrt{n}$ shows: for any nonzero effect, no matter how trivial, collecting enough data will eventually make it ``statistically significant''---which is why significance is not the same as importance.
\end{toystory}

The $t$-statistic for testing $H_0: \beta_j = 0$ is
\begin{equation}\label{eq:tstat}
  t = \frac{\bhat_j}{\se(\bhat_j)}.
\end{equation}
Since $\se(\bhat_j) = O(1/\sqrt{n})$, for any fixed $\beta_j \neq 0$:
\[
  |t| \;\to\; \infty \quad \text{as } n \to \infty.
\]

\begin{keyinsight}
With enough data, \emph{every} non-zero effect becomes statistically significant, no matter how small.
With $n = 50{,}000$ observations (NovaMart at store-month level), customer satisfaction has $p < 0.001$ despite a coefficient of \$0.003 per satisfaction point---economically meaningless.
\end{keyinsight}

\subsection{Statistical Power}

\begin{definition}[Power]
The \textbf{power} of a test is the probability of correctly rejecting $H_0$ when $H_0$ is false:
\[
  \text{Power} = P\!\left(|t| > c_\alpha \;\middle|\; \beta_j \neq 0\right),
\]
where $c_\alpha$ is the critical value at significance level~$\alpha$.
\end{definition}

\begin{proposition}[Power Increases with $n$ and Effect Size]
For a two-sided $t$-test at level $\alpha$ with true coefficient $\beta_j$ and standard deviation $\sigma$,
\[
  \text{Power} \approx \Phi\!\left(\frac{|\beta_j|}{\sigma/\sqrt{n}} - z_{\alpha/2}\right),
\]
where $\Phi$ is the standard normal CDF and $z_{\alpha/2}$ is the upper $\alpha/2$ quantile.
Power increases with both $|\beta_j|$ (effect size) and $\sqrt{n}$.
\end{proposition}

\subsection{The Multiple Testing Problem}

Running $m$ independent hypothesis tests at level $\alpha$ produces
\[
  \E[\text{false positives}] = m \cdot \alpha.
\]
With $m = 1{,}000{,}000$ tests at $\alpha = 0.05$, we expect 50{,}000 false positives.

\begin{repairbox}
The \textbf{Bonferroni correction} tests each hypothesis at level $\alpha/m$, controlling the family-wise error rate at~$\alpha$.
For genome-wide association studies ($m \approx 10^6$), the corrected threshold is $\alpha/m = 5 \times 10^{-8}$.
Always report effect sizes and confidence intervals alongside $p$-values.
\end{repairbox}

\begin{numermark}[title={The Keyboard That Saves 0.1 Seconds}]
A new keyboard layout saves an average of $\beta = 0.1$ seconds per email ($\sigma = 2$ seconds). Is this ``significant''?

\begin{center}
\begin{tabular}{@{}cccc@{}}
\toprule
Sample size $n$ & $\se = \sigma/\sqrt{n}$ & $t = 0.1/\se$ & Significant? \\
\midrule
5 & 0.894 & 0.11 & No \\
50 & 0.283 & 0.35 & No \\
5{,}000 & 0.028 & 3.54 & Yes ($p < 0.001$) \\
\bottomrule
\end{tabular}
\end{center}

The effect never changed---it was always 0.1 seconds. Only the sample size grew. With 5{,}000 emails, saving one tenth of a second becomes ``highly significant'' ($p < 0.001$). Statistical significance tells you ``the effect is not zero.'' It says nothing about whether the effect matters.
\end{numermark}

\begin{vizspec}[title={``The Magnifying Glass That Can't Measure'' --- Sample Size Simulator}]
\textbf{Try it:} A slider controls sample size $n$ (50 to 100{,}000). Three linked panels show the $|t|$-statistic rising, the $p$-value plummeting, and the confidence interval narrowing---all while the effect size stays fixed at a tiny value.

\textbf{The aha moment:} The dashboard goes from gray ``not significant'' to screaming red ``HIGHLY SIGNIFICANT ($p < 0.001$)'' while the actual effect label never changes. Significance is a property of $n$, not of importance.
\end{vizspec}

% ============================================================
% 6. SPECIFICATION ERROR (NB 4)
% ============================================================
\section{Why Your Model Is Wrong: Specification Error}\label{sec:specification}

\subsection{The Problem of Functional Form}

\begin{toystory}[title={The Straight Ruler on a Winding Road}]
You're driving up a mountain on a road that curves back and forth. A friend asks you to describe the road, but you only have a straight ruler. You lay it on the map across the section you've driven so far, and it gives a decent summary---``on average, the road goes uphill at about this angle.'' But then you extend that straight line past where you've actually been, and it says you should be 500 feet in the air by now, flying off a cliff. The ruler was the \emph{best} straight summary of the curvy road you had data for, but using it to predict beyond that range was a disaster.

This is exactly what the linear projection under misspecification gives you: OLS finds the best-fitting straight line through curved reality, which can look reasonable in-sample but diverges catastrophically when you extrapolate.
\end{toystory}

Suppose the true relationship is $Y = f(X) + \varepsilon$ for some nonlinear function $f$, but we fit the linear model $Y = X\beta + \varepsilon$.

\begin{theorem}[Linear Projection Under Misspecification]\label{thm:linproj}
If $Y = f(X) + \varepsilon$ with $\E[\varepsilon \mid X] = 0$, and we estimate $Y = X\beta + u$ by OLS, then the OLS estimator converges to
\[
  \beta^* = \underset{\beta}{\arg\min}\; \E\!\left[(Y - X\beta)^2\right],
\]
which is the best linear approximation to $f(X)$ in the mean-squared-error sense.
This $\beta^*$ equals the true $\beta$ only if $f$ is linear.
\end{theorem}

In the NovaMart DGP, $f(X_1) = \beta_1 \log(X_1)$.
A linear model in $X_1$ provides a reasonable approximation within the observed data range but diverges catastrophically when extrapolating.

\subsection{The Extrapolation Danger}

If the true relationship is $Y = 8\log(X_1) + \cdots$ and we fit $Y = \tilde\beta_1 X_1 + \cdots$:
\begin{itemize}[nosep]
  \item At $X_1 = 100$ (within sample): moderate overestimation.
  \item At $X_1 = 2{,}000$ (extrapolation): linear prediction overshoots by a factor of $\approx$3.
  \item The confidence band widens but does \emph{not} contain the truth.
\end{itemize}

\begin{repairbox}
Fit the correctly specified model by transforming the predictor:
\[
  Y = \beta_0 + \beta_1 \log(X_1) + \beta_2 X_2 + \varepsilon.
\]
Extrapolation accuracy improves dramatically.
However, even the log model is a local approximation---at extreme extrapolation, all parametric models eventually fail.
\end{repairbox}

\begin{numermark}[title={Square Plots and Fence Lengths}]
The true relationship between the area of a square plot ($X$, in m$^2$) and its side length ($Y = \sqrt{X}$, in m) is nonlinear. We fit a straight line to five plots:

\begin{center}
\begin{tabular}{@{}cccc@{}}
\toprule
Plot & $X$ (area) & True $Y = \sqrt{X}$ & Fitted $\hat{Y} = 1.24 + 0.16X$ \\
\midrule
A & 1 & 1.00 & 1.40 \\
B & 4 & 2.00 & 1.88 \\
C & 9 & 3.00 & 2.68 \\
D & 16 & 4.00 & 3.80 \\
E & 25 & 5.00 & 5.25 \\
\bottomrule
\end{tabular}
\end{center}

Maximum in-sample error: just $0.40$. Looks fine! Now extrapolate:
\begin{itemize}[nosep]
  \item At $X = 100$: line predicts $\hat{Y} = 17.3$. Truth: $\sqrt{100} = 10$. Off by 73\%.
  \item At $X = 10{,}000$: line predicts $\hat{Y} = 1{,}606$. Truth: $\sqrt{10{,}000} = 100$. Off by a factor of 16.
\end{itemize}
The best linear approximation is only ``best'' within the range where you asked the question.
\end{numermark}

\begin{vizspec}[title={``The Comfort Zone'' --- Extrapolation Explorer}]
\textbf{Try it:} Drag a prediction point along the $x$-axis. Within the data range, the background is calm blue and the error bar is tiny. Drag past the data boundary: the background turns progressively red, the error bar explodes, and the warning escalates from ``Mild extrapolation'' to ``CATASTROPHIC: prediction is fiction.''

\textbf{Toggle:} Show the correctly specified $\log$ model alongside the line. The log model tracks truth even far out---the contrast is stark.
\end{vizspec}

% ============================================================
% 7. OVERFITTING AND BIAS-VARIANCE (NB 5)
% ============================================================
\section{Why Your $R^2$ Is Lying: Overfitting}\label{sec:overfitting}

\subsection{The Bias--Variance Decomposition}

\begin{toystory}[title={The Three Dartboard Players}]
Three friends play darts. The first always throws at the same wrong spot---consistently two inches left of the bullseye. Every throw lands in the same cluster, just in the wrong place. That's high bias, low variance. The second friend aims perfectly and averages a bullseye, but her throws scatter wildly all over the board. That's low bias, high variance. The third friend accepts that the dartboard sits on a wobbly wall, so even a perfect throw gets jostled randomly. That wobble is irreducible noise---nobody can beat it. The best player is the one who aims at the right spot AND throws consistently, finding the sweet spot between the first two friends' mistakes.

This is exactly what the bias-variance decomposition says: total prediction error equals systematic mis-aim (bias squared) plus scatter in your throws (variance) plus the wobble you can never control (irreducible noise).
\end{toystory}

\begin{theorem}[Bias--Variance Decomposition]\label{thm:biasvar}
For any estimator $\hat{f}$ of the regression function $f$, the expected prediction error at a point $x_0$ decomposes as
\begin{equation}\label{eq:biasvar}
  \E\!\left[(Y - \hat{f}(x_0))^2\right]
  = \underbrace{\left(f(x_0) - \E[\hat{f}(x_0)]\right)^2}_{\text{Bias}^2}
  + \underbrace{\Var(\hat{f}(x_0))}_{\text{Variance}}
  + \underbrace{\sigma^2_\varepsilon}_{\text{Irreducible noise}}.
\end{equation}
\end{theorem}

\begin{proof}
Let $\mu = \E[\hat{f}(x_0)]$.  Then
\begin{align*}
  \E[(Y - \hat{f})^2]
  &= \E[(f + \varepsilon - \hat{f})^2] \\
  &= \E[(f - \hat{f})^2] + \E[\varepsilon^2] + 2\E[\varepsilon(f - \hat{f})] \\
  &= \E[(f - \hat{f})^2] + \sigma^2_\varepsilon,
\end{align*}
where the cross-term vanishes because $\varepsilon$ is independent of $\hat{f}$ (trained on separate data) and has mean zero.
Now expand:
\[
  \E[(f - \hat{f})^2]
  = (f - \mu)^2 + \E[(\hat{f} - \mu)^2]
  = \Bias^2(\hat{f}) + \Var(\hat{f}). \qedhere
\]
\end{proof}

\subsection{Overfitting in Practice}

Adding polynomial terms to a regression always increases training $R^2$ (since $R^2$ is monotonically non-decreasing in the number of regressors).
In the NovaMart demo:
\begin{center}
\begin{tabular}{@{}ccc@{}}
\toprule
\textbf{Degree} & \textbf{Training $R^2$} & \textbf{Test $R^2$} \\
\midrule
1 & 0.71 & 0.69 \\
3 & 0.82 & 0.78 \\
7 & 0.91 & 0.63 \\
15 & 0.97 & $<0$ \\
\bottomrule
\end{tabular}
\end{center}
A degree-15 polynomial achieves $R^2 = 0.97$ on training data but predicts \emph{worse than the sample mean} on held-out data.

\subsection{Cross-Validation}

\begin{definition}[$k$-Fold Cross-Validation]
Partition the data into $k$ equally sized folds.  For each fold $j$, train the model on all data except fold~$j$ and compute the prediction error on fold~$j$.
The CV estimate of prediction error is
\[
  \text{CV}(k) = \frac{1}{k} \sum_{j=1}^{k} \MSE_j.
\]
\end{definition}

\begin{repairbox}
Using cross-validation, the log-linear model (which matches the true DGP) achieves the lowest CV error, outperforming all polynomial models.
The right model family matters more than the degree of flexibility.
\end{repairbox}

\begin{warningbox}
Cross-validation itself can overfit if the number of candidate model specifications searched is large.
The best CV $R^2$ inflates with the breadth of the search.
\end{warningbox}

\begin{numermark}[title={Three Models, Three Training Sets}]
True relationship: $Y = 2X$. We predict at $X = 6$ (true $Y = 12$) using three training sets of $(X, Y)$ pairs and three models:

\begin{center}
\begin{tabular}{@{}lccc@{}}
\toprule
& \textbf{Model A} (constant) & \textbf{Model B} (line) & \textbf{Model C} (degree-4) \\
\midrule
Prediction (Set 1) & 7 & 10.8 & 8.0 \\
Prediction (Set 2) & 7 & 13.2 & 16.0 \\
Prediction (Set 3) & 7 & 12.2 & $-13.0$ \\
\midrule
Average prediction & 7.0 & 12.1 & 3.7 \\
Bias$^2$ & 25.0 & $\approx 0$ & 69.4 \\
Variance & 0.0 & 1.0 & 149.6 \\
\textbf{Total MSE} & \textbf{25.0} & \textbf{1.0} & \textbf{219.0} \\
\bottomrule
\end{tabular}
\end{center}

Model C has zero training error on every dataset, yet its predictions swing from $+16$ to $-13$. The line (Model B) wins decisively. Perfect in-sample fit amplifies noise into wild out-of-sample predictions.
\end{numermark}

\begin{vizspec}[title={``The U-Shaped Trap'' --- Bias-Variance Explorer}]
\textbf{Try it:} A slider controls polynomial degree (1 to 15). The left panel shows a stacked area chart of Bias$^2$ + Variance + Noise forming the classic U-shaped MSE curve. The right panel overlays 20 fitted curves from different training sets.

\textbf{The aha moment:} At degree 1, twenty nearly identical lines all miss the true curve (stable but wrong). At degree 15, twenty polynomials pass through their training points perfectly but look completely different from each other---wild oscillations everywhere (unstable and absurd). Press ``Animate'' and watch the sweet spot at degree 3--4.
\end{vizspec}

% ============================================================
% 8. BAD CONTROLS AND CAUSATION (NB 6)
% ============================================================
\section{Why Your Causal Claim Is Wrong: Endogeneity and Bad Controls}\label{sec:causation}

\subsection{Endogeneity}

\begin{toystory}[title={Why Every Movie Star's Spouse Seems Boring}]
Among famous actors, their spouses seem to be unusually un-funny or un-attractive. Does fame make you marry dull people? No---you're only looking at people who are already famous. People become famous either by being extremely talented or by being extremely good-looking (or both). Once you filter for ``is famous,'' the talented-but-plain people and the gorgeous-but-untalented people get mixed together, creating a fake tradeoff that doesn't exist in the general population. By restricting your attention to a group defined by the outcome of the very thing you're studying, you manufactured a relationship from thin air.

This is exactly what collider bias describes: when you control for a variable that is \emph{caused by} the things you're studying, you create a spurious association that didn't exist before---more controls made the answer worse, not better.
\end{toystory}

\begin{definition}[Endogeneity]
A regressor $X_j$ is \emph{endogenous} if $\Cov(X_j, \varepsilon) \neq 0$.
This can arise from omitted variables, measurement error, or simultaneous causation (feedback loops).
\end{definition}

In the NovaMart DGP, revenue and ad spend are mutually related through demand~$U$.
OLS captures the correlation but cannot distinguish causal direction.

\subsection{Collider Bias from Bad Controls}

The causal graph for NovaMart is:
\[
  U \longrightarrow X_1 \longrightarrow Y \longrightarrow X_3, \qquad U \longrightarrow Y.
\]
Including $X_3$ (satisfaction) as a regressor conditions on a descendant of $Y$, which opens a spurious association path through $Y$.

\begin{proposition}[Collider Bias]
If $X_3 = \alpha Y + \eta_3$ (i.e., $X_3$ is a post-treatment variable), then the regression of $Y$ on $(X_1, X_3)$ yields a biased estimate of the causal effect of $X_1$ on~$Y$, even if $X_1$ is otherwise exogenous.
The bias arises because conditioning on a descendant of $Y$ partially conditions on $Y$ itself.
\end{proposition}

\begin{warningbox}
The advice ``control for everything available'' is wrong.
Controlling for post-treatment variables or colliders \emph{introduces} bias rather than removing it.
Variable selection for causal inference requires understanding the causal structure, not just statistical association.
\end{warningbox}

\begin{numermark}[title={Talent, Luck, and the Fame Filter}]
People have Talent (0 or 1) and Luck (0 or 1), independently assigned. A person is Famous if they have Talent OR Luck (or both).

\begin{center}
\begin{tabular}{@{}cccc@{}}
\toprule
Type & Talent & Luck & Famous? \\
\midrule
I & 0 & 0 & No \\
II & 0 & 1 & Yes \\
III & 1 & 0 & Yes \\
IV & 1 & 1 & Yes \\
\bottomrule
\end{tabular}
\end{center}

\textbf{Full population:} $P(\text{Talent}=1) = 1/2$, $P(\text{Luck}=1) = 1/2$, and $P(\text{both}) = 1/4 = 1/2 \times 1/2$. Talent and Luck are \textbf{independent}. Correlation $= 0$.

\textbf{Among the Famous only} (drop Type I): If someone is talented (Types III, IV), $P(\text{Lucky}) = 1/2$. If not talented (Type II only), $P(\text{Lucky}) = 1$. Suddenly, knowing a famous person is talented makes them \emph{less} likely to be lucky!

The observation ``famous people are either talented OR lucky but rarely both'' is not evidence that the universe is fair---it is a mathematical artifact of selecting on a collider.
\end{numermark}

\begin{vizspec}[title={``The Paradox of Conditioning'' --- Collider Explorer}]
\textbf{Try it:} A scatterplot of Talent vs.\ Luck for 1{,}000 people. All dots are gray; no correlation visible. Toggle ``Condition on Fame'' and watch: most dots fade, the highlighted subset reveals a clear negative slope. A DAG panel alongside animates: a box appears around the Fame node, and a dashed red ``spurious association'' arc draws between Talent and Luck.

\textbf{Slider:} Fame threshold---from 10th to 90th percentile. At low thresholds (most people famous), the spurious correlation is weak. At high thresholds (top 5\%), it becomes extreme.
\end{vizspec}

\subsection{Instrumental Variables}

\begin{definition}[Instrument]
A variable $Z$ is a valid \emph{instrument} for $X$ in the equation $Y = X\beta + \varepsilon$ if:
\begin{enumerate}[nosep]
  \item \textbf{Relevance:} $\Cov(Z, X) \neq 0$.
  \item \textbf{Exclusion:} $\Cov(Z, \varepsilon) = 0$ (the instrument affects $Y$ only through $X$).
\end{enumerate}
\end{definition}

\begin{theorem}[IV Consistency]\label{thm:iv}
If $Z$ is a valid instrument, the IV estimator
\begin{equation}\label{eq:iv}
  \biv = (Z^\top X)^{-1} Z^\top Y
\end{equation}
is consistent for~$\beta$:
\[
  \plim\, \biv
  = \plim \frac{Z^\top Y / n}{Z^\top X / n}
  = \frac{\Cov(Z, Y)}{\Cov(Z, X)}
  = \beta.
\]
\end{theorem}

In the NovaMart context, if ad prices varied randomly across regions (a natural experiment), the regional price variation serves as an instrument for ad spend.

\begin{warningbox}
Weak instruments ($\Cov(Z, X) \approx 0$) cause the IV estimator to be highly variable and unreliable.
When the first-stage $F$-statistic is below 10, IV confidence intervals are extremely wide and the estimator may be more biased than OLS.
\end{warningbox}

% ============================================================
% 9. CONFRONTING REAL DATA (NB 7)
% ============================================================
\section{The Real World: When All Pathologies Strike at Once}\label{sec:realworld}

\subsection{The Lalonde Dataset}

Notebook~7 transitions from simulated data to the Lalonde (1986) dataset, a benchmark in causal inference.
The question: does a job training programme increase earnings?

\begin{itemize}[nosep]
  \item \textbf{Observational estimate (OLS):} negative coefficient---training appears to \emph{reduce} earnings.
  \item \textbf{Experimental estimate (RCT):} positive and large ($\approx$\$1{,}800).
  \item The two estimates do not even have the same sign.
\end{itemize}

All standard diagnostics pass: VIFs are low, Breusch--Pagan is non-significant, robust SEs barely change.
The problem is not in the model---it is in the \emph{data}.
The observational and experimental groups differ systematically on unobservables.

\subsection{Sensitivity Analysis: Cinelli--Hazlett Framework}

\begin{definition}[Robustness Value]
The \textbf{robustness value} (RV) is the minimum strength of association (measured in partial~$R^2$) that an unobserved confounder would need to have with both the treatment and the outcome to explain away the observed result.
\end{definition}

Formally, define
\[
  \text{RV}_{q,\alpha} = \min\{r : \text{adjusted estimate crosses } q \cdot \bhat \text{ at level } \alpha\}.
\]
A small RV indicates fragility: a weak confounder suffices to nullify the result.
A large RV indicates robustness.

\begin{keyinsight}
The Lalonde observational estimate has a very small robustness value: the result is fragile.
A confounder with even modest partial $R^2$ with both treatment and outcome would explain away the finding.
This is why the observational and experimental estimates disagree.
\end{keyinsight}

% ============================================================
% 10. REGRESSION DISCONTINUITY DESIGN (NB 8)
% ============================================================
\section{The Redemption: Regression Discontinuity Design}\label{sec:rdd}

\subsection{When Design Earns Trust}

Regression works brilliantly when the world provides a natural experiment.
The \textbf{regression discontinuity design} (RDD) exploits a sharp threshold to identify causal effects.

\begin{definition}[Sharp RDD]
Suppose treatment is assigned deterministically by a \emph{running variable} $R$ at a cutoff~$c$:
\[
  D_i = \mathbf{1}(R_i \geq c).
\]
Under continuity of potential outcomes at~$c$, the \textbf{local average treatment effect} at the cutoff is
\begin{equation}\label{eq:rdd}
  \tau_{\text{RDD}} = \lim_{r \downarrow c}\, \E[Y \mid R = r] - \lim_{r \uparrow c}\, \E[Y \mid R = r].
\end{equation}
\end{definition}

\begin{theorem}[RDD Identification]\label{thm:rdd}
If $\E[Y_0 \mid R = r]$ and $\E[Y_1 \mid R = r]$ are continuous at $r = c$, then
\[
  \tau_{\text{RDD}} = \E[Y_1 - Y_0 \mid R = c].
\]
The design eliminates confounding \emph{locally} at the cutoff because units just above and just below the threshold are nearly identical in all characteristics.
\end{theorem}

\subsection{Example: Scholarship Cutoff}

Students with test scores just above a threshold receive scholarships; those just below do not.
Near the cutoff, the two groups are essentially identical.
The jump in outcomes at the threshold identifies the causal effect of the scholarship.

\subsection{Bandwidth Choice}

The RDD estimator faces a bias--variance tradeoff controlled by the \emph{bandwidth}~$h$:
\begin{itemize}[nosep]
  \item \textbf{Narrow bandwidth:} low bias (observations are close to the cutoff), high variance (few observations).
  \item \textbf{Wide bandwidth:} low variance (more observations), high bias (observations far from the cutoff may not be comparable).
\end{itemize}

Optimal bandwidth selection (e.g., Imbens--Kalyanaraman or Calonico--Cattaneo--Titiunik) balances these two forces.

\begin{repairbox}
The RDD does not require controlling for confounders.
The design itself ensures comparability near the cutoff.
When the autopsy report is run on the RDD estimate, all indicators are green: the regression is trustworthy because the \emph{design} is trustworthy.
\end{repairbox}

% ============================================================
% 11. THE FRISCH-WAUGH-LOVELL THEOREM (NB 1.5)
% ============================================================
\section{What ``Controlling For'' Actually Means: The Frisch--Waugh--Lovell Theorem}\label{sec:fwl}

\subsection{Motivation}

When we say ``the coefficient on $X_1$, controlling for $X_2$,'' we do \emph{not} mean ``holding $X_2$ constant.''
We mean: the coefficient from regressing the part of $Y$ unexplained by $X_2$ on the part of $X_1$ unexplained by $X_2$.

\begin{theorem}[Frisch--Waugh--Lovell]\label{thm:fwl}
Consider the regression $Y = X_1 \beta_1 + X_2 \beta_2 + \varepsilon$.
Let $M_2 = I - X_2(X_2^\top X_2)^{-1}X_2^\top$ be the residual-maker matrix that projects orthogonally to the column space of~$X_2$.
Then
\begin{equation}\label{eq:fwl}
  \bhat_1 = (X_1^\top M_2 X_1)^{-1} X_1^\top M_2 Y.
\end{equation}
That is, $\bhat_1$ from the full regression is \emph{identical} (to numerical precision) to the coefficient obtained by:
\begin{enumerate}[nosep]
  \item Regressing $Y$ on $X_2$ and taking residuals: $\tilde{Y} = M_2 Y$.
  \item Regressing $X_1$ on $X_2$ and taking residuals: $\tilde{X}_1 = M_2 X_1$.
  \item Regressing $\tilde{Y}$ on $\tilde{X}_1$ (no intercept).
\end{enumerate}
\end{theorem}

\begin{proof}
The OLS normal equations for the partitioned model $Y = X_1\beta_1 + X_2\beta_2 + \varepsilon$ give
\[
  \begin{pmatrix} X_1^\top X_1 & X_1^\top X_2 \\ X_2^\top X_1 & X_2^\top X_2 \end{pmatrix}
  \begin{pmatrix} \bhat_1 \\ \bhat_2 \end{pmatrix}
  =
  \begin{pmatrix} X_1^\top Y \\ X_2^\top Y \end{pmatrix}.
\]
From the second block, $\bhat_2 = (X_2^\top X_2)^{-1} X_2^\top (Y - X_1 \bhat_1)$.
Substituting into the first block and simplifying:
\begin{align*}
  X_1^\top X_1 \bhat_1 + X_1^\top X_2 (X_2^\top X_2)^{-1} X_2^\top(Y - X_1\bhat_1) &= X_1^\top Y \\
  X_1^\top\!\big[I - X_2(X_2^\top X_2)^{-1} X_2^\top\big] X_1\, \bhat_1 &= X_1^\top\!\big[I - X_2(X_2^\top X_2)^{-1} X_2^\top\big] Y \\
  X_1^\top M_2 X_1\, \bhat_1 &= X_1^\top M_2 Y.
\end{align*}
Hence $\bhat_1 = (X_1^\top M_2 X_1)^{-1} X_1^\top M_2 Y$.
\end{proof}

\begin{keyinsight}
``Controlling for'' means residualising---stripping out the linear relationship with the control variables from \emph{both} sides of the equation.
The partial regression plot (also called an added-variable plot) visualises this: it plots $\tilde{Y}$ against $\tilde{X}_1$.
The slope of this scatterplot \emph{is} $\bhat_1$ from the full multiple regression.
\end{keyinsight}

% ============================================================
% 12. DIAGNOSTICS TOOLKIT SUMMARY
% ============================================================
\section{Diagnostic Toolkit Summary}\label{sec:diagnostics}

The following diagnostics are used across all notebooks:

\begin{center}
\begin{tabular}{@{}llp{6.5cm}@{}}
\toprule
\textbf{Diagnostic} & \textbf{Tests For} & \textbf{Interpretation} \\
\midrule
Residuals vs.\ Fitted & Nonlinearity, heteroscedasticity & Residuals should show no pattern; funnel shape indicates heteroscedasticity \\
Normal Q--Q plot & Non-normality of residuals & Points should lie on the 45$^\circ$ line \\
Scale--Location & Heteroscedasticity & $\sqrt{|\text{standardised residuals}|}$ should be flat against fitted values \\
Residuals vs.\ Leverage & Influential observations & Points with high leverage \emph{and} large residuals (Cook's $d > 1$) are influential \\
\midrule
VIF & Multicollinearity & $\text{VIF}_j > 10$ suggests problematic collinearity \\
Breusch--Pagan & Heteroscedasticity & $H_0$: homoscedasticity; reject if $p < \alpha$ \\
Durbin--Watson & Autocorrelation & Statistic near 2 indicates no autocorrelation \\
Jarque--Bera & Non-normality & $H_0$: residuals are normally distributed \\
\bottomrule
\end{tabular}
\end{center}

\begin{warningbox}
Passing all diagnostics does not guarantee a valid causal estimate.
In the Lalonde example (Notebook~7), every diagnostic passes, yet the OLS estimate has the wrong sign.
Diagnostics check the model's internal consistency, not whether the research design is sound.
\end{warningbox}

% ============================================================
% 13. THEOREM MAP
% ============================================================
\section{Theorem Map: How the Results Connect}\label{sec:theoremmap}

The theorems in this course form a logical chain where each generalises the previous:

\begin{center}
\renewcommand{\arraystretch}{1.3}
\begin{tabular}{@{}clp{7.5cm}@{}}
\toprule
\textbf{NB} & \textbf{Theorem} & \textbf{Generalisation Relationship} \\
\midrule
1 & Omitted Variable Bias & Foundation: bias from omitting a relevant variable \\
2 & Sandwich Variance & Relaxes $\Var(\varepsilon) = \sigma^2 I$; OVB is a special case of endogeneity \\
3 & Power Analysis & $t = \bhat/\se(\bhat)$; builds on correct SEs from NB~2 \\
4 & Linear Projection & Best linear approx.\ to $f(X)$; relaxes $f$ linear \\
5 & Bias--Variance Tradeoff & $\MSE = \Bias^2 + \Var + \sigma^2$; generalises specification error \\
6 & IV Consistency & If $\E[\varepsilon X] \neq 0$ but $\E[\varepsilon Z] = 0$; generalises OVB repair \\
7 & Sensitivity Analysis & Robustness value quantifies fragility of any estimate \\
8 & RDD Identification & Continuity at cutoff $\Rightarrow$ causal identification \\
1.5 & Frisch--Waugh--Lovell & What ``controlling for'' means algebraically \\
\bottomrule
\end{tabular}
\end{center}

The meta-lesson: each theorem relaxes an assumption that justified the previous conclusion.
Classical variance is a special case of the sandwich estimator.
The linear model is a special case of all models.
OVB is a special case of endogeneity.

% ============================================================
% CONCLUSION
% ============================================================
\section{Conclusion}

These notes have formalised the eight core failure modes of OLS regression demonstrated in the Regression Autopsy lab series:

\begin{enumerate}[nosep]
  \item Omitted variable bias shifts your coefficient.
  \item Heteroscedasticity corrupts your standard errors.
  \item Large samples make trivial effects ``significant.''
  \item Misspecification makes extrapolation fiction.
  \item Overfitting makes memorisation look like learning.
  \item Bad controls and endogeneity destroy causal claims.
  \item Real data confronts you with all failures simultaneously.
  \item Sound research design (RDD) earns regression its trust back.
\end{enumerate}

The unifying lesson is not that regression should be abandoned, but that its outputs must be accompanied by an honest assessment of the assumptions under which they are valid.
The diagnostics tell you whether the \emph{model} is internally consistent; only the \emph{research design} can tell you whether the estimate is credible.

\vspace{1cm}
\begin{center}
\textit{This doesn't mean you should stop using regression.\\
It means you're finally ready to start.}
\end{center}

\end{document}
